\section{Proof-Batching Tools}\label{sec:tools}

This section presents the two batching tools used repeatedly in the protocol.
The first handles heterogeneous linear identities whose natural domains
differ. The second merges many same-form local nonlinear checks into one
larger virtual instance. Together they cover the two batching patterns that
recur throughout the rest of the construction.
Standard divisibility-based zero checks and same-domain batching already appear
in prior proof systems such as Aurora and CQ++ \cite{aurora19,cq++24}. Our
contribution here is to isolate two reusable tools: domain-lifting batching for
heterogeneous linear claims, and interleaving batching for aligned same-form
nonlinear claims.

\subsection{Domain-Lifting Batching}

As explained earlier, most linear constraints in our construction are first
expressed as zero-checks over some evaluation domain $\mathbb{H}_n$. The
following theorem shows how to batch many such claims together even when their
natural domains differ. The key idea is to lift each claim from its original
domain to a common larger domain, so the lifted claim is still a zero-check
but now over the same domain as the others. After lifting, we can apply an
ordinary same-domain batch to combine all claims into one.

\begin{theorem}[Domain-lifting batching]
\label{thm:domain-lifting-batching}
Let
\[
2 \leq n_1 \leq n_2 \leq \cdots \leq n_m = N
\]
be powers of two. For each $i \in [m]$, let $F_i(X)$ be a virtual polynomial
and suppose the target claim is
\[
F_i(X)=Q_i(X)\nu_{n_i}(X)
\]
for some polynomial $Q_i(X)$. Define the lifted polynomial
\[
\widehat{F}_i(X):=
F_i(X)\frac{\nu_N(X)}{\nu_{n_i}(X)}.
\]
Then:
\begin{enumerate}
  \item $F_i(X)=Q_i(X)\nu_{n_i}(X)$ holds if and only if
  $\widehat{F}_i(X)$ is divisible by $\nu_N(X)$.
  \item For a verifier challenge $\eta$, all $m$ claims can be combined into
  the single identity
  \begin{equation}\label{eq:domain-lifting-batching}
  \sum_{i=1}^{m}\eta^{i-1}\widehat{F}_i(X)=Q(X)\nu_N(X).
  \end{equation}
  \item If at least one original claim is false, then
  \eqref{eq:domain-lifting-batching} holds for at most $m-1$ values of $\eta$, i.e. the soundness error is at most $(m-1)/|\mathbb{F}|$.
\end{enumerate}
\end{theorem}

\begin{proof}
The first item is immediate from
\[
F_i(X)=Q_i(X)\nu_{n_i}(X)
\iff
F_i(X)\frac{\nu_N(X)}{\nu_{n_i}(X)} = Q_i(X)\nu_N(X).
\]
For soundness, once the lifted residuals are reduced modulo $\nu_N(X)$, the
left-hand side of \eqref{eq:domain-lifting-batching} becomes a polynomial in $\eta$
of degree at most $m-1$. If at least one original claim is false, this
polynomial is nonzero, so a uniformly random $\eta$ causes accidental
acceptance with probability at most $(m-1)/|\mathbb{F}|$.
\end{proof}

Theorem~\ref{thm:domain-lifting-batching} is the protocol's main tool for batching
heterogeneous linear identities. Later we apply it to routing equations,
padding-induced tensor equalities, and other zero-check-style relations whose
natural domains are not identical.

\subsection{Interleaving Batching}

The intended setting for interleaving batching is that many checks share the
same local rule and the same public parameters, but arrive as separate 
PCS and witness on possibly different power-of-two domains. We first record the alignment lemma
used to place smaller objects on larger domains.

\begin{lemma}[Row replication]
\label{lem:row-replication}
Let $\mathbf{v}\in\mathbb{F}^m$ and let $v(X)$ be its encoding over
$\mathbb{H}_m$. Let $\mathbf{V}\in\mathbb{F}^{r\times m}$ be the matrix whose
every row equals $\mathbf{v}$. Then the encoding polynomial of $\mathbf{V}$ over
$\mathbb{H}_{rm}$ is
\[
V(X)=v(X^r).
\]
\end{lemma}

\begin{proof}
For every $i\in[r]$ and $j\in[m]$,
\[
V(\omega_{rm}^{im+j})
= \mathbf{V}[i,j]
= \mathbf{v}[j]
= v(\omega_m^j)
= v((\omega_{rm}^{im+j})^r),
\]
where the last equality uses $\omega_{rm}^r=\omega_m$. Since both sides have
degree less than $rm$ and agree on $\mathbb{H}_{rm}$, they are identical.
\end{proof}

\begin{definition}[Pairwise interleaving]
\label{def:pairwise-interleaving}
Let $f(X)$ and $g(X)$ be the encoding polynomials of two length-$n$ vectors
over $\mathbb{H}_n$, where $n$ is a power of two. Define
\[
\operatorname{even}_n(X):=\frac{X^n+1}{2},
\qquad
\operatorname{odd}_n(X):=\frac{1-X^n}{2}.
\]
Their evaluations on $\mathbb{H}_{2n}$ are the alternating selectors
$(1,0,1,0,\ldots)$ and $(0,1,0,1,\ldots)$. The pairwise interleaving of
$f(X)$ and $g(X)$ is defined as the virtual polynomial
\[
\operatorname{Int}_n(f,g)(X)
:=
\operatorname{even}_n(X)\,f(X)
+
\operatorname{odd}_n(X)\,g(\omega_{2n}^{-1}X).
\]
Its evaluation vector over $\mathbb{H}_{2n}$ is
\[
(f[0],g[0],f[1],g[1],\ldots,f[n-1],g[n-1]).
\]
\end{definition}

\begin{theorem}[Interleaving batching]
\label{thm:interleaving-batch}
For each $i\in[m]$, let $u_i(X)$ denote the polynomial encoding a vector of
length $n_i$ over $\mathbb{H}_{n_i}$, where each $n_i$ is a power of two. For
some (non-linear) constraint $\varphi$, suppose the target claim is that all
$m$ pointwise relations
\[
\forall i \in [m],\ \forall a\in\mathbb{H}_{n_i}\quad
\varphi\!\left(
u_i(a)
\right)=0.
\]
There exists an algorithm that combines these $m$ claims into one
check over a single virtual polynomial $\widehat{u}(X)$ over domain
$\mathbb{H}_{\widehat{n}}$, using only domain lifts and pairwise interleaving
and introducing no new PCS commitment. Using two scans over power-of-two
domain sizes, it runs in $O(m+\log \hat{n})$ time. This algorithm guarantees:
\begin{enumerate}
  \item all original pointwise claims hold if and only if
  \[
  \forall z\in\mathbb{H}_{\widehat{n}}\quad
  \varphi\!\left(
  \widehat{u}(z)
  \right)=0;
  \]
  \item the final domain size satisfies
  \[
  \widehat{n}
  =
  2^{\left\lceil \log_2(n_1+\cdots+n_m)\right\rceil}.
  \]
  \item any opening claim for $\widehat{u}(X)$ is reduced to one opening for each
  original $u_i(X)$, and the verifier can check the claim in
  $O(m\log \widehat{n})$ time.
\end{enumerate}
\end{theorem}

\begin{proof}
Start from the active multiset $\{(u_i(X),n_i)\}_{i\in[m]}$.
We first store all pairs with same domain size in the same bucket.
Then we perform two scans over the buckets.
In the first scan, process the buckets from small to large domain size and
merge equal-sized pairs whenever a bucket contains at least two active pairs.
Each such merge uses pairwise interleaving as in
\cref{def:pairwise-interleaving}. After this pass, every bucket contains at
most one active pair, and the total size $\sum_i n_i$ is unchanged.

In the second scan, repeatedly take the two smallest active pairs. If their
sizes differ, lift the smaller one to the larger domain using
$X\mapsto X^r$ with $r=n_j/n_i$; then interleave the pair. Repeat
until one active pair $(\widehat{u}(X),\widehat{n})$ remains.

Because the domain sizes are powers of two, the two scans touch only
$O(\log \hat{n})$ buckets, and each merge or lift removes one active pair, so the
implementation runs in $O(m+\log \hat{n})$ time.

For property~(1), lifting is exactly the setting of \cref{lem:row-replication}:
if $r=n_j/n_i$, then $X\mapsto X^r$ repeats the entries of $u_i(X)$ on the
larger domain $\mathbb{H}_{n_j}$. So the pointwise relation $\varphi(\cdot)=0$
holds before lifting if and only if it holds after lifting. Pairwise
interleaving preserves the same relation for the same reason,
since all entries are preserved.
Repeating these two facts over the two scans proves property~(1).

For property~(2), note that the first scan preserves the total size
$S:=n_1+\cdots+n_m$. If we denote the
sizes after the first scan by $n_1'<\cdots<n_s'$, then
\[
\sum_{i=1}^m n_i = \sum_{i=1}^s n_s'\implies n_s' = 2^{\left\lfloor \log_2(n_1+\cdots+n_m)\right\rfloor}.
\]
If $n_s' = 2^{\log_2(n_1+\cdots+n_m)}$,
there $s=1$ and there is no second scan, and the property~(2) holds.
Otherwise,
the second scan pushes the size to $2 n_s'$,
and thus
\[
\widehat{n}= 2 n_s' = 2^{\left\lceil \log_2(n_1+\cdots+n_m)\right\rceil}.
\]

For property~(3), note that for interleaving,
the verifier can compute the sparse selectors by itself.
And for row-replication $u(X)\mapsto u(X^r) = v(X)$, the opening of $v(x)$ for any $x\in \mathbb{F}$ is trivially reduced to the opening of $u(x^r)$.
Hence, the verifier can check any opening claim for $\widehat{u}(X)$ by checking one opening claim for each original $u_i(X)$,
which incurs $O(m\log \widehat{n})$ time to compute the selectors and check the claims.
\end{proof}

\subsection{Explicit Lookup Batching Example}

As a concrete application for interleaving batching, suppose we have several lookup families indexed by
$i\in[s]$. For each family, there is a public table $\mathbf{t}_i$ and a
collection of vector queries $\mathbf{f}_{i,1},\ldots,\mathbf{f}_{i,m_i}$ such
that every query satisfies $\mathbf{f}_{i,j}\subseteq \mathbf{t}_i$.

We batch them in three steps. First, for each fixed $i$, we interleave the
queries $\mathbf{f}_{i,1},\ldots,\mathbf{f}_{i,m_i}$ into one family-level
query vector $\mathbf{f}_i$. Second, we shift each table by a random challenge
$\alpha$ and form $\widetilde{\mathbf{t}}_i:=\alpha^i+\mathbf{t}_i$, so that
different tables are separated with high probability. We then interleave the
shifted tables into one combined table vector $\mathbf{t}$. Third, we
interleave the family-level query vectors $\mathbf{f}_1,\ldots,\mathbf{f}_s$
into one global query vector $\mathbf{f}$.

This yields a single lookup claim
\[
\mathbf{f}\subseteq \mathbf{t},
\]
which can be proved with one lookup argument instead of one argument per
family. The benefit is that the prover amortizes the lookup proof across all
tables while keeping the query structure explicit. The random shifts prevent
cross-table collisions, and the interleaving steps preserve the original
within-family lookup relations.

A second example is fixed-point quantization. Suppose several checks all prove
relations of the form
\[
a_i = \left\lfloor \frac{b_i}{2^q} \right\rfloor
\]
for the same public divisor $2^q$ and the same rounding rule. These are again
same-form local checks: after aligning the domains of the witness vectors, we
interleave them and obtain one larger quantization check. This is the pattern
behind the fixed-point gain, leaf-value, and normalization relations that
appear later in the construction.

Theorem~\ref{thm:interleaving-batch} is useful when the protocol produces many
small nonlinear checks of the same form, and lookup batching across different
tables is one of the main examples in our construction.
